\section*{Chapter \ref{ch:in:bank-quant} --- Bank Quant}

\begin{solution}{exo:in:bank-quant:1}
From \$128\,300 to \$158\,100: 23.2\% in dollars. Prices rose by $135.8312/114.325=1.188$, so the real rise is
$1.232/1.188-1=3.7\%$.
\end{solution}

\begin{solution}{exo:in:bank-quant:2}
Arrangement (a): the two units report to different members of senior management (the chief risk officer and the head of
markets), and the validator is not under the developer or the owner.
\end{solution}

\begin{solution}{exo:in:bank-quant:3}
Tier one: a full validation every year (400 hours), no separate review, and 10\% of 400 for material changes: 440 hours.
Tier three: $60/3+8\times\frac23+0.1\times60=31.3$ hours.
\end{solution}

\begin{solution}{exo:in:bank-quant:4}
With a one-year cycle for tier two, each tier-two model costs $160+16=176$ hours instead of 104: the inventory needs 82.1
validators instead of 63.9.
\end{solution}

\begin{solution}{exo:in:bank-quant:5}
The change term doubles, from 13\,498 to 26\,996 hours: 72.4 validators.
\end{solution}

\begin{solution}{exo:in:bank-quant:6}
The strategist has a stake in the model being approved, knows what it was built to do rather than what it fails to do,
and would review assumptions he chose. Every rule of the chapter (the 2011 US guidance's ``stake'', the UK statement's
separate reporting lines, the ECB's arrangements) excludes it; the strategist's knowledge enters as documentation and
tests that the validator reviews.
\end{solution}

\begin{solution}{exo:in:bank-quant:7}
Each draw tiers 1\,000 random score triples and applies the method: the mean is 64.0 validators and the 5th--95th
percentile range 60.1--67.9. The randomness of the inventory moves the answer by about three validators either way;
the choice of hours moves it far more.
\end{solution}

\begin{solution}{exo:in:bank-quant:8}
The validators report to the head of quantitative analytics, who owns the library they review: they are under the
developer, so \texttt{independence} returns \texttt{fail}, whatever the team's name. Arrangement (b) requires two units
that meet only at a member of senior management.
\end{solution}

\begin{solution}{pb:in:bank-quant:1}
\begin{enumerate}
\item As in the chapter's four definitions: the desk's model and tool builder; the builder of the shared pricing library;
the builder of risk and capital models; the independent reviewer of others' models.
\item Goldman Sachs: desk strats ``sit on trading floors and provide value for both internal \& external clients through
cutting-edge models, predictive analytics, and high-quality trading tools''.
\item So that the same trade has one price in the trading system, the risk engine and the accounts.
\item Many risk and capital models need a supervisor's approval; a material change needs it again.
\item Between the desk and risk: it simulates exposures like a \omterm{def:in:bank-quant:risk}{risk quant}, prices charges like a strategist, and its
numbers go into the accounts.
\item Validators should have no stake in the result; reporting lines may support independence, but it is judged by
actions and outcomes; each model reviewed at least annually.
\item The 2026 revised interagency guidance (SR~26-2; OCC Bulletin 2026-13): risk-based, with the timing and frequency of
validation varying with the model.
\item Separate reporting lines for validators and for developers and owners, and tiering by materiality and complexity.
\item (a) Different members of senior management; (b) the same member, different units; (c) separate staff in one unit.
\item Internal audit evaluates whether model risk management is rigorous and effective; it does not duplicate development
or validation.
\item $H=\sum_t n_t(F_t/k_t+R_t(1-1/k_t)+cF_t)$ hours; $H/h$ validators.
\item Twice the materiality score plus complexity and uncertainty: 10 or more is tier one (yearly), 7 to 9 tier two
(every two years), less tier three (every three years).
\item 102, 403 and 495.
\item 43.9\%, for 10.2\% of the models: they are revalidated in full every year at the highest hours.
\item 18.1 validators (82.1 against 63.9).
\item 63.9 validators; about 2.6 more for each percentage point of tier-one share.
\item Illustrative: the score probabilities, the hours, the change rate and the productive hours. From rules: annual
review in the 2011 US guidance, annual validation of internal models in the ECB guide, tiering in the UK statement; the
cycles by tier are Book~6's illustrative policy.
\item Arrangement (a), the most robust.
\item Internal audit, which checks that the function works; and the supervisors, who review the bank's framework.
\item At the bank's current tiering the inventory needs about 64 validators, and each percentage point of models moved
into tier one adds about three more. The team reports to the chief risk officer, separately from the developers and
owners, which meets every rule the chapter cites.
\end{enumerate}
\end{solution}

\begin{solution}{iq:in:bank-quant:1}
Compare with an independent implementation over a grid (strikes, maturities, volatilities, rates), check limits (zero
volatility, zero maturity, deep in and out of the money), put--call parity, Greeks against finite differences, and
behaviour at extreme inputs.
\iqlookfor{an independent benchmark, limits and identities, not one test price.}
\end{solution}

\begin{solution}{iq:in:bank-quant:2}
Find out why before changing anything: the same inputs (volatility surface, barrier monitoring, rebate, dividends)? The
model (local or stochastic volatility, which price barriers differently)? A difference in the market's price is evidence
about the model to be understood, not a price to be copied.
\iqlookfor{separating inputs from model risk.}
\end{solution}

\begin{solution}{iq:in:bank-quant:3}
Two and a half are expected; seven or more has probability 1.4\% under a correct model, so it is evidence against it, and
it falls in the Basel yellow zone (five to nine). Check first for data and position errors and for clustering in a
stress period, then the model.
\iqlookfor{a binomial test and a look at the causes.}
\end{solution}

\begin{solution}{iq:in:bank-quant:4}
Measure the change on every desk's positions before release, explain it, have it validated as a model change, agree a date
with the desks and product control, and release with the old version available for comparison.
\iqlookfor{impact analysis and change control.}
\end{solution}

\begin{solution}{iq:in:bank-quant:5}
Its intended use and the decisions it feeds; its limitations as stated by the developer; the data and assumptions; and what
testing was done against what benchmark. A document that does not say what the model must not be used for is a finding.
\iqlookfor{use, limitations and evidence.}
\end{solution}

\begin{solution}{iq:in:bank-quant:6}
The difference is model risk that today's surface cannot resolve, because a cliquet depends on forward volatility.
Price with one, reserve or report the difference as model uncertainty, try a model with stochastic volatility, and give
the validators both calibrations and the difference.
\iqlookfor{forward-volatility dependence and a model reserve.}
\end{solution}
